\section{Introduction}
\label{sec:introduction}

Vision-Language Models (VLMs) have demonstrated remarkable capabilities in understanding and generating content at the intersection of visual and textual modalities. However, their deployment in high-stakes domains is severely hindered by the phenomenon of \textit{hallucination}—where models generate plausible but visually ungrounded or factually incorrect information. Recent studies indicate that even state-of-the-art models suffer from hallucination rates as high as 30\% in complex reasoning tasks, manifesting as object fabrication, attribute misidentification, and reasoning failures.

\subsection{Motivation and Problem Statement}
The challenge of mitigating hallucinations in VLMs is compounded by the "black-box" nature of large pre-trained models. Existing approaches typically fall into two categories: training-based and training-free methods. Training-based methods, such as RLHF (Reinforcement Learning from Human Feedback) or DPO (Direct Preference Optimization), require expensive computational resources and extensive annotated datasets. While effective, they can sometimes degrade the model's general capabilities and may require additional adaptation to comfortably extend to entirely new domains without retraining.

Conversely, training-free methods, such as decoding strategies or post-hoc correction, offer a more flexible alternative. However, current training-free solutions face notable limitations:
\begin{enumerate}
    \item \textbf{Separation of Detection and Mitigation:} Many frameworks treat detection and correction as separate phases, often detecting errors only after the full response is generated, which precludes real-time intervention.
    \item \textbf{Static Knowledge Bases:} Approaches relying on Retrieval-Augmented Generation (RAG) typically use static knowledge graphs (KGs). These static structures cannot adapt to novel entities or the evolving context of a dynamic visual scene, leading to "knowledge obsolescence."
    \item \textbf{Unidirectional Verification:} Most verification mechanisms focus solely on text-to-image consistency, neglecting the bidirectional nature of cross-modal understanding. This results in systems that may correctly identify objects but fail to verify their relationships or attributes against external factual knowledge.
    \item \textbf{Perception-Reasoning Gap:} A significant portion of hallucinations stems not from perceptual errors but from flawed reasoning. Current methods often focus on object-level detection, missing subtle logical inconsistencies or temporal errors in complex scenes.
\end{enumerate}

\subsection{Proposed Solution: CMVKG-Guard}
To address these gaps, we propose \textbf{CMVKG-Guard} (Cross-Modal Verified Knowledge Graph Guard), a novel training-free framework that integrates real-time hallucination detection and correction. unlike existing pipeline approaches, CMVKG-Guard operates as an intelligent middleware, dynamically constructing a Hierarchical Multimodal Knowledge Graph (DH-MMKG) during inference. This graph serves as a structured grounding source, evolving with the conversation and the visual context.

Our system employs a Cross-Modal Verification Engine (CMVE) that validates generated tokens across three layers: (1) \textit{Visual-Textual Alignment} to ensure fidelity to the input image; (2) \textit{Knowledge Grounding} to verify facts against the constructed MMKG and external knowledge bases; and (3) \textit{Reasoning Chain Verification} to ensure logical consistency. Furthermore, we introduce an Adaptive Confidence Calibration (ACC) mechanism that dynamically adjusts detection thresholds based on query complexity, significantly reducing false positives.

\subsection{Contributions}
The key contributions of this work are as follows:
\begin{itemize}
    \item \textbf{Simultaneous token-level detection and correction.} Unlike prior methods that treat detection and correction as sequential offline phases, CMVKG-Guard intercepts generation token-by-token in real time, correcting errors before they cascade. This \textit{online} verification paradigm is a departure from all existing training-free approaches.
    \item \textbf{Dynamic Hierarchical Multimodal Knowledge Graph (DH-MMKG).} We introduce a novel per-inference graph that self-constructs from the input image and evolves with the dialogue context. This removes the ``knowledge obsolescence'' limitation of static RAG approaches, as every entity in the visual scene is guaranteed to be represented with up-to-date relational structure.
    \item \textbf{Three-layer Unified Verification Score (UVS).} We formalise a non-differentiable, heuristic verification engine (CMVE) that fuses visual, factual, and logical consistency signals into a single score. The algorithmic contribution is the design of this multi-axis check that addresses three orthogonal hallucination causes (perceptual, knowledge, reasoning) simultaneously, unlike single-axis baselines.
    \item \textbf{Step-aware Adaptive Calibration.} We contribute a dynamic threshold function that penalises later generation steps more heavily, explicitly countering the ``snowballing'' effect of cascading hallucinations in long outputs — a problem not addressed by any prior threshold-based method.
    \item \textbf{Empirical validation.} Extensive experiments on POPE, CHAIR, and MMHalBench against eight baselines (including recent NeurIPS/ICML/ICLR 2025 methods) show that CMVKG-Guard reduces object hallucinations by 84\% on CHAIR and achieves 91.5\% F1 on POPE, with a real-time latency overhead of 12ms per token.
\end{itemize}
